% Auto-extracted from ../tfpt-42.tex for the TFPT 4.5 clean split.
% Source body assigned by the TFPT 4.5 clean split.

% ---- Carrier algebraic normal form through abelian index: source lines 1306-2314 ----
\subsection{Carrier algebraic normal form}
\label{sec:minimal-carrier-proof}

\begin{lemma}[Algebraic normal form from exterior signatures]
\label{thm:carrier-minimality-k}
\label{thm:carrier-minimality-kpp}
Let
\[
E=E_-\oplus E_+
\]
be a finite chiral carrier with two simple factors of ranks
\[
b:=\dim E_-,
\qquad
s:=\dim E_+,
\]
and let
\[
S^+=\Lambda^{\mathrm{even}}E
\]
be the positive half-spinor packet. Let
\[
X=q_-P_-+q_+P_+
\]
be a primitive integer two-point generator with
\[
q_-<0<q_+,
\qquad
bq_-+sq_+=0.
\]
Assume only
\[
\dim \Lambda^2 E_+=1,
\qquad
\Lambda^2 E_- \cong E_-^\vee.
\]
Then necessarily
\[
s=2,
\qquad
b=3,
\qquad
g:=b+s=5.
\]
Hence
\[
\dim \Lambda^{\mathrm{even}}E=16.
\]
The even exterior algebra decomposes into exactly the six sectors
\[
(1,1)_0,\quad
(3,2)_{1/6},\quad
(\bar 3,1)_{-2/3},\quad
(\bar 3,1)_{1/3},\quad
(1,2)_{-1/2},\quad
(1,1)_1,
\]
and the primitive trace condition therefore has the unique solution
\[
(q_-,q_+)=(-2,3).
\]
Consequently
\[
X=-2P_3+3P_2,
\qquad
Y:=\frac{X}{6}
=
\mathrm{diag}\!\left(-\frac13,-\frac13,-\frac13,\frac12,\frac12\right).
\]
\end{lemma}

\begin{proof}
From
\[
\dim\Lambda^2 E_+=\binom{s}{2}=1
\]
one gets $s=2$. For the color factor one always has
\[
\Lambda^{b-1}E_- \cong E_-^\vee.
\]
The assumption
\[
\Lambda^2E_- \cong E_-^\vee
\]
therefore forces $b-1=2$, hence $b=3$. Thus $g=5$.

Now
\[
S^+=\Lambda^0E\oplus\Lambda^2E\oplus\Lambda^4E,
\]
and for $E=E_3\oplus E_2$,
\[
\Lambda^2E
=
\Lambda^2E_3
\oplus
(E_3\otimes E_2)
\oplus
\Lambda^2E_2,
\]
\[
\Lambda^4E
=
(\Lambda^3E_3\otimes E_2)
\oplus
(\Lambda^2E_3\otimes \Lambda^2E_2).
\]
These are exactly the six sectors
\[
(1,1)_0,\ (3,2)_{1/6},\ (\bar 3,1)_{-2/3},\ (\bar 3,1)_{1/3},\ (1,2)_{-1/2},\ (1,1)_1.
\]
Hence
\[
\dim S^+=2^{5-1}=16.
\]
Finally
\[
3q_-+2q_+=0
\]
with primitive integers and $q_-<0<q_+$ has the unique solution
\[
(q_-,q_+)=(-2,3).
\]
\end{proof}

\begin{corollary}[One chiral family emerges]
Under the hypotheses of \TFPTxref{thm:carrier-minimality-k}, the positive half-spinor packet $S^+$
carries exactly one chiral Standard Model family including $\nu^c$.
\end{corollary}

\begin{proof}
The six-sector decomposition displayed in \TFPTxref{thm:carrier-minimality-k} is exactly
\[
(1,1)_0\oplus(3,2)_{1/6}\oplus(\bar 3,1)_{-2/3}\oplus(\bar 3,1)_{1/3}\oplus(1,2)_{-1/2}\oplus(1,1)_1,
\]
namely one chiral Standard Model family together with $\nu^c$.
\end{proof}

\begin{remark}[Logical status of the carrier normal form]
The lemma above is the algebraic normal form of the former open carrier target Theorem~K. It is a rigidity
statement under two exterior signatures:
\[
\dim \Lambda^2 E_+=1,
\qquad
\Lambda^2 E_- \cong E_-^\vee.
\]
Within those signatures the packet size
\[
\dim\Lambda^{\mathrm{even}}E=16
\]
and the one-family decomposition are outputs, while the occupancy identities move behind the
carrier closure as derived consistency statements. On the retained branch, the compact Higgs index
first fixes the bosonic rank and \TFPTxref{thm:yukawa-forces-32,cor:yukawa-discharges-kpp} then
discharge these carrier-signature premises internally from the actual branch Yukawa tensor.
\end{remark}

\begin{corollary}[Carrier norm]
On the minimal carrier one has
\[
Y
=
\mathrm{diag}\!\left(-\frac13,-\frac13,-\frac13,\frac12,\frac12\right),
\qquad
\gamma=\tr_EY^2=\frac56.
\]
\end{corollary}

Hence from now on
\[
E=E_3\oplus E_2,
\qquad
X=-2P_3+3P_2,
\qquad
P_3=\frac{\mathbf 1-\epscar}{2},
\qquad
P_2=\frac{\mathbf 1+\epscar}{2},
\]
and the hypercharge operator is the normalized primitive two-point generator
\begin{equation}
\label{eq:explicit-hypercharge}
Y
=
\mathrm{diag}\!\left(-\frac13,-\frac13,-\frac13,\frac12,\frac12\right)
=
\frac{X}{6}
=
-\frac13 P_3+\frac12 P_2
=
\frac1{12}\mathbf 1+\frac5{12}\epscar.
\end{equation}

\begin{lemma}[Charge-generator / carrier-polarization inversion]
On the rigid carrier, the carrier polarization and the primitive two-point generator are recovered by
\begin{align}
\epscar&=\frac{2X-\mathbf 1}{5}=\frac{12Y-\mathbf 1}{5},
\label{eq:hypercharge-seam-inversion}\\
\qquad
X&=\frac{5\epscar+\mathbf 1}{2},
\nonumber\\
(2X-\mathbf 1)^2&=25\,\mathbf 1,
\nonumber\\
X^2-X-6\,\mathbf 1&=0,
\nonumber\\
6Y^2-Y-\mathbf 1&=0.
\label{eq:carrier-polynomial}
\end{align}
Consequently the carrier projectors are
\[
P_2=\frac{\mathbf 1+\epscar}{2}=\frac{2(3Y+\mathbf 1)}{5},
\qquad
P_3=\frac{\mathbf 1-\epscar}{2}=\frac{3(\mathbf 1-2Y)}{5}.
\]
\end{lemma}

\begin{proof}
The affine relations between $(X,Y)$ and $\epscar$ are inverted algebraically. Squaring the first
relation and using $\epscar^2=\mathbf 1$ gives the quadratic generator polynomial
$X^2-X-6\,\mathbf 1=0$, and dividing by $36$ yields the hypercharge polynomial. The projector
formulas then follow from $P_2=(\mathbf 1+\epscar)/2$ and $P_3=(\mathbf 1-\epscar)/2$.
\end{proof}

\begin{lemma}[Primitive trace-balanced carrier]
\label{lem:primitive-trace-balanced-carrier}
Let $Y$ be an endomorphism of a finite-dimensional carrier $E$ such that
\[
6Y^2-Y-\mathbf 1=0,
\qquad
\tr_EY=0,
\]
and assume the realization is primitive in the sense that $E$ is not a nontrivial direct sum of
identical trace-balanced $Y$-modules. Then
\[
\operatorname{Spec}(Y)=\left\{-\frac13,\frac12\right\},
\qquad
\dim\ker\!\left(Y+\frac13\mathbf 1\right)=3,
\qquad
\dim\ker\!\left(Y-\frac12\mathbf 1\right)=2.
\]
Equivalently,
\[
\rank E=5,
\qquad
E=E_3\oplus E_2,
\qquad
Y=-\frac13P_3+\frac12P_2.
\]
\end{lemma}

\begin{proof}
The quadratic polynomial factors as
\[
(2Y-\mathbf 1)(3Y+\mathbf 1)=0,
\]
so the only eigenvalues are
\[
y_+=\frac12,
\qquad
y_-=-\frac13.
\]
Write
\[
m_+:=\dim\ker(Y-y_+\mathbf 1),
\qquad
m_-:=\dim\ker(Y-y_-\mathbf 1).
\]
Then
\[
0=\tr_EY=\frac12\,m_+-\frac13\,m_-,
\]
so
\[
3m_+=2m_-.
\]
Because $\gcd(2,3)=1$, every positive integer solution has the form
\[
(m_+,m_-)=(2k,3k).
\]
Primitivity forces $k=1$, hence $(m_+,m_-)=(2,3)$. This is exactly the rigid $3+2$ split, with
the displayed spectral form of $Y$.
\end{proof}

\begin{remark}[Compressed carrier decoder]
The carrier can therefore be entered from the shorter algebraic normal form
\[
6Y^2-Y-\mathbf 1=0,
\qquad
\tr_EY=0,
\qquad
\text{primitive realization}
\Longrightarrow
E=E_3\oplus E_2.
\]
The exterior-signature route of \TFPTxref{thm:carrier-minimality-k} remains the theorem-level discharge
of these algebraic hypotheses from the actual branch geometry.
\end{remark}

The hard carrier package is summarized visually in \Cref{fig:minimal-carrier-glance}. The point is
that the same minimal data are readable in four equivalent ways: geometrically as an oriented
double cover with deck reflection and boundary polarization, algebraically as the trace-balanced
quadratic carrier decoder, combinatorially as the rigid $3+2$ split, and spectrally as the
primitive two-point generator $X$ with normalized hypercharge $Y=X/6$.

\begin{figure}[t]
\centering
\footnotesize
\begin{tikzpicture}[
    scale=0.88,
    transform shape,
    >=Latex,
    slot3/.style={draw=blue!60!black, fill=blue!6, rounded corners, minimum width=0.72cm, minimum height=0.72cm, align=center},
    slot2/.style={draw=orange!70!black, fill=orange!9, rounded corners, minimum width=0.72cm, minimum height=0.72cm, align=center},
    info/.style={draw=black!55, rounded corners, fill=gray!3, align=center, inner sep=5pt}
]
\node[font=\bfseries] at (-3.5,1.95) {oriented double cover};
\draw[draw=gray!70!black, fill=gray!6, rounded corners] (-5.4,0.55) rectangle (-1.7,1.35);
\draw[draw=gray!70!black, fill=gray!6, rounded corners] (-5.4,-1.35) rectangle (-1.7,-0.55);
\node at (-3.55,0.95) {upper sheet};
\node at (-3.55,-0.95) {lower sheet};
\draw[very thick, purple!70!black] (-5.6,0) -- (-1.45,0);
\node[anchor=west, purple!70!black] at (-1.38,0) {seam $\Sigma$};
\draw[<->, thick] (-3.55,0.42) -- (-3.55,-0.42);
\node[anchor=west] at (-3.4,0) {$\iota$};

\node[font=\bfseries] at (2.35,1.95) {rigid carrier split};
\node[slot3] (c1) at (0.55,0.95) {$c_1$};
\node[slot3, right=2mm of c1] (c2) {$c_2$};
\node[slot3, right=2mm of c2] (c3) {$c_3$};
\node[slot2, right=6mm of c3] (w1) {$w_1$};
\node[slot2, right=2mm of w1] (w2) {$w_2$};
\node[align=center, blue!60!black] at ($(c2.south)+(0,-0.45)$) {$E_3=P_3E$\\three color slots};
\node[align=center, orange!75!black] at ($(w1.south)+(0.45,-0.45)$) {$E_2=P_2E$\\two weak slots};
\draw[decorate, decoration={brace, amplitude=4pt}] ($(c1.north west)+(-0.07,0.12)$) -- ($(c3.north east)+(0.07,0.12)$);
\node[blue!60!black] at ($(c2.north)+(0,0.42)$) {$3$};
\draw[decorate, decoration={brace, amplitude=4pt}] ($(w1.north west)+(-0.07,0.12)$) -- ($(w2.north east)+(0.07,0.12)$);
\node[orange!75!black] at ($(w1.north east)+(0.2,0.42)$) {$2$};

\node[info] at (-0.45,-2.3) {%
    $\displaystyle
    X=-2P_3+3P_2,
    \qquad
    Y=\frac{X}{6}
    =-\frac13 P_3+\frac12 P_2
    =\frac1{12}\mathbf 1+\frac5{12}\iota
    \qquad
    \iota=\frac{2X-\mathbf 1}{5}
    $\\[2pt]
    $\displaystyle
    P_3=\frac{\mathbf 1-\iota}{2},
    \qquad
    P_2=\frac{\mathbf 1+\iota}{2},
    \qquad
    \mathrm{Spec}(Y)=\left\{-\frac13,\frac12\right\}
    $};
\end{tikzpicture}
\caption{Rigid carrier at a glance. The seam involution, the primitive two-point generator, and the
normalized hypercharge operator carry the same hard information: geometrically they separate the
two sheets, and algebraically they separate the $3+2$ carrier split through the projectors $P_3$
and $P_2$.}
\label{fig:minimal-carrier-glance}
\end{figure}

\begin{remark}[Two-point carrier algebra]
Because the minimal carrier satisfies $6Y^2-Y-\mathbf 1=0$, every analytic function of $Y$ reduces to
\[
f(Y)=A_f\,\mathbf 1+B_f\,Y,
\qquad
A_f=\frac{2f(1/2)+3f(-1/3)}{5},
\qquad
B_f=\frac{6\bigl(f(1/2)-f(-1/3)\bigr)}{5}.
\]
Thus the hard carrier is effectively a two-point algebra supported on the spectral values $-1/3$ and $1/2$, and the projectors $P_3,P_2$ are already polynomial functions of $Y$.
\end{remark}
